\begin{table}[htbp]
\centering
\footnotesize
\caption{Problema G (openai/gpt-4o-mini): LLM versus Perceptron treinado nos mesmos exemplos, acurácia contra o rótulo real.}
\label{tab:bloco3_G_faseE_LLM_vs_perceptron_gpt4mini}
\begin{tabular}{lllrrrrrrrrrr}
\toprule
 &  &  & \multicolumn{5}{c}{Acur.\ LLM vs real (\%)} & \multicolumn{5}{c}{Acur.\ Perceptron vs real (\%)} \\ \cmidrule(lr){4-8} \cmidrule(lr){9-13}
Subproblema & Classes & Atributos & n=0 & n=4 & n=10 & n=20 & n=40 & n=0 & n=4 & n=10 & n=20 & n=40 \\
\midrule
G1 (normalizado) & A / B & peso/altura & 25,6 & 19,6 & 16,7 & 12,2 & 13,3 & — & 16,2 & 14,4 & 12,6 & 11,1 \\
G1 (normalizado) & A / B & x1/x2 & 30,0 & 20,7 & 21,1 & 16,3 & 15,9 & — & 22,4 & 13,3 & 14,4 & 10,7 \\
\textbf{G2 (kg e m)} & \textbf{Mulher / Homem} & \textbf{peso\_kg/altura\_m} & \textbf{88,9} & \textbf{83,7} & \textbf{82,2} & \textbf{81,5} & \textbf{78,9} & \textbf{—} & \textbf{83,9} & \textbf{88,1} & \textbf{86,3} & \textbf{83,3} \\
G2 (kg e m) & Mulher / Homem & x1/x2 & 62,2 & 62,2 & 69,3 & 71,1 & 68,5 & — & 85,0 & 83,3 & 80,7 & 78,5 \\
G2 (kg e m) & A / B & peso\_kg/altura\_m & 40,0 & 51,5 & 48,9 & 43,7 & 40,4 & — & 32,2 & 38,9 & 34,4 & 31,1 \\
\bottomrule
\multicolumn{13}{p{0.95\linewidth}}{\footnotesize Fase E nas 30 amostras de teste; o especialista é a melhor métrica da Fase A externa. Perceptron = Perceptron diagonal treinado nos mesmos n exemplos, no mesmo espaço R$^n$ da melhor métrica (R$^2$ quando a eleita tem 2 atributos; se a Fase A externa colapsou, sem melhor métrica, a Fase E de G não roda e a linha não existe). Acurácia (vs rótulo real) do LLM e do Perceptron; médias das sementes e sorteios. Nas linhas com classes neutras A/B a acurácia literal do LLM pode sair invertida (ver bloco3\_table\_02).} \\
\end{tabular}
\end{table}
